\documentclass[11pt,a4paper]{article}
\usepackage{fontspec,amsmath,amssymb,amsthm,microtype}
\setmainfont{lmroman10-regular.otf}[BoldFont=lmroman10-bold.otf,ItalicFont=lmroman10-italic.otf,BoldItalicFont=lmroman10-bolditalic.otf]
\usepackage{xcolor}
\usepackage[most]{tcolorbox}
\usepackage{mathrsfs}
\providecommand{\lt}{<}
\providecommand{\gt}{>}
\setmonofont{DejaVuSansMono.ttf}[BoldFont=DejaVuSansMono-Bold.ttf,ItalicFont=DejaVuSansMono-Oblique.ttf,BoldItalicFont=DejaVuSansMono-BoldOblique.ttf,Scale=0.86]
\newfontfamily\leanfallback{STIXTwoMath-Regular.otf}[Scale=0.86]
\definecolor{leanrule}{gray}{0.62}
\definecolor{leanback}{gray}{0.962}
\definecolor{leantitle}{gray}{0.22}
\definecolor{leanmuted}{gray}{0.38}
\newtcolorbox{leanbox}[1]{enhanced,breakable,colback=leanback,colframe=leanrule,
  boxrule=0.4pt,arc=0.6mm,left=2.4mm,right=2.4mm,top=1.8mm,bottom=1.8mm,
  before skip=0.8em,after skip=1em,fonttitle=\small\bfseries,coltitle=leantitle,
  colbacktitle=leanback,titlerule=0.3pt,title={#1}}
\newlength{\leancharwidth}
\newenvironment{leancode}{\par\smallskip\begingroup\ttfamily\footnotesize\raggedright
  \settowidth{\leancharwidth}{M}\setlength{\parskip}{0pt}\setlength{\parindent}{0pt}
  \frenchspacing}{\par\endgroup\smallskip}
\newcommand{\leanline}[2]{\hangindent=\dimexpr#1\leancharwidth+4\leancharwidth\relax
  \hangafter=1\noindent\hspace*{#1\leancharwidth}#2\par}
\newcommand{\leansource}[1]{\par{\footnotesize\color{leanmuted}#1}\par}
\newcommand{\leanlabel}[1]{\par\smallskip\noindent{\small\bfseries #1}\par}
\usepackage[margin=25mm]{geometry}
\usepackage[unicode,hidelinks]{hyperref}
\hypersetup{pdfauthor={LangevinLean Project}}
\setlength{\emergencystretch}{3em}
\setlength{\parskip}{0.4em}
\setlength{\parindent}{0pt}
\allowdisplaybreaks[2]
\raggedbottom
\title{\texorpdfstring{A step length and iteration count for the KLMC error tolerance}{A step length and iteration count for the KLMC error tolerance}}
\author{LangevinLean Project}
\date{}
\begin{document}\maketitle
Draft mathematical exposition; these notes may contain mistakes.


The convergence theorem becomes an explicit choice of parameters: choose a positive step, an integer iteration count, and obtain a error at most \(\varepsilon\). The result concerns the constructed sampler. Neither a recurrence nor a local-error estimate is left as a premise.



The first-order sampling theorem is converted into a positive-step, integer-iteration prescription. Explicit bias allocation and a rounded logarithmic transient bound handle finite initialization and requested accuracy.


\section{\texorpdfstring{Verify that the geometric factor is admissible}{Verify that the geometric factor is admissible}}


Set \(\alpha=3mh/(4\gamma)\). The curvature, friction and step assumptions imply \(0<\alpha\le1\). The proof reuses the exponential-coefficient estimate: \(e^{-mh/\gamma}+(3/4)Mh^2\le1-\alpha\). Its left side is positive, so the right side is nonnegative. This justifies every later geometric and logarithmic estimate.



A positive contraction rate ensures the logarithmic threshold is meaningful. The upper bound on the rate ensures the geometric base is nonnegative. Both properties are consequences of the friction and step restrictions, and are needed before solving the error inequality for the iteration number.

\begin{leanbox}{Lean: \texttt{klmc\_\allowbreak{}rate\_\allowbreak{}fraction\_\allowbreak{}le\_\allowbreak{}one}}
\leansource{\texttt{klmc\_\allowbreak{}rate\_\allowbreak{}fraction\_\allowbreak{}le\_\allowbreak{}one} \quad KineticLangevin/\allowbreak{}Accuracy.lean, lines 10–20}
\begin{leancode}
\leanline{0}{theorem\ klmc\_rate\_fraction\_le\_one\ (m\ M\ γ\ h\ :\ ℝ)\ (hm\ :\ 0<m)\ (hmM\ :\ m≤M)}
\leanline{4}{(hγ\ :\ 0<γ)\ (hfriction\ :\ m+M≤γ\textasciicircum{}2)\ (hh\ :\ 0≤h)\ (hstep\ :\ 4*γ*M*h≤m)\ :}
\leanline{4}{3*m*h/(4*γ)≤1}
\end{leancode}
\leanlabel{How the Lean proof follows the argument}
\texttt{klmc\_\allowbreak{}rate\_\allowbreak{}fraction\_\allowbreak{}le\_\allowbreak{}one} gives the formal counterpart of this step. The main intermediate facts are \texttt{hb}, \texttt{he}, \texttt{h\allowbreak{}M}, \texttt{hp}, \texttt{hid} in that order. The final step is the algebraic normalization or ordered arithmetic displayed in the proof.
\end{leanbox}

\section{\texorpdfstring{Apply the scalar error allocation to the law}{Apply the scalar error allocation to the law}}


The KLMC theorem gives \(W_2(\nu_n,\pi)\le\sqrt2(1-\alpha)^nD+b\). The scalar error bound allocates \(b\le0.94\varepsilon\) and \((1-\alpha)^nD\le\varepsilon/24\). The sum is at most \(\varepsilon\). The threshold is the natural ceiling of \(\max\{0,\log(24D/\varepsilon)/\alpha\}\); if \(D=0\), no initial-error iterations are needed. The finite-transport hypothesis prevents the real-valued Wasserstein convention from hiding an infinite distance.



When the initial discrepancy is already below its transient allocation, the maximum-logarithm prescription selects zero iterations. Otherwise the exponential upper bound for the geometric factor is below the allocated error after the selected ceiling count. The bias allocation then completes the total-error bound.

\begin{leanbox}{Lean: \texttt{klmc\_\allowbreak{}accuracy\_\allowbreak{}of\_\allowbreak{}error\_\allowbreak{}budgets}}
\leansource{\texttt{klmc\_\allowbreak{}accuracy\_\allowbreak{}of\_\allowbreak{}error\_\allowbreak{}budgets} \quad KineticLangevin/\allowbreak{}Accuracy.lean, lines 23–38}
\begin{leancode}
\leanline{0}{theorem\ klmc\_accuracy\_of\_error\_budgets\ \{d\ :\ ℕ\}\ \{μ\ :\ Measure\ (CanonicalCoupling.H\ d)\}}
\leanline{4}{[IsProbabilityMeasure\ μ]\ (f\ :\ CanonicalCoupling.H\ d\ →\ ℝ)\ (m\ :\ ℝ)\ (M\ γ\ :\ ℝ≥0)}
\leanline{4}{(hf\ :\ Differentiable\ ℝ\ f)\ (hsc\ :\ StrongConvexOn\ univ\ m\ f)\ (hLip\ :\ LipschitzWith\ M\ (gradient\ f))}
\leanline{4}{(hmM\ :\ m≤(M:ℝ))\ (hm\ :\ 0<m)\ (hγ\ :\ 0<(γ:ℝ))\ (hfriction\ :\ m+(M:ℝ)≤(γ:ℝ)\textasciicircum{}2)}
\leanline{4}{(h\ :\ ℝ≥0)\ (hh\ :\ 0<(h:ℝ))\ (hstep\ :\ 4*(γ:ℝ)*(M:ℝ)*(h:ℝ)≤m)}
\leanline{4}{(hfinite\ :\ ULA.wasserstein2\ μ\ (ULA.gibbsMeasure\ f)≠⊤)\ (ε\ :\ ℝ)\ (hε\ :\ 0<ε)}
\leanline{4}{(hbias\ :\ (M:ℝ)*(h:ℝ)*Real.sqrt\ (2*(d:ℝ))/m≤(47/50:ℝ)*ε)\ (n\ :\ ℕ)}
\leanline{4}{(hn\ :\ ULA.wasserstein2Real\ μ\ (ULA.gibbsMeasure\ f)=0\ ∨}
\leanline{6}{firstOrderIterationCount\ (3*m*(h:ℝ)/(4*(γ:ℝ)))}
\leanline{8}{(ULA.wasserstein2Real\ μ\ (ULA.gibbsMeasure\ f))\ ε≤n)\ :}
\leanline{4}{ULA.wasserstein2Real\ (CanonicalCoupling.klmcLaw\ f\ γ\ h\ μ\ n)\ (ULA.gibbsMeasure\ f)≤ε}
\end{leancode}
\leanlabel{How the Lean proof follows the argument}
\texttt{klmc\_\allowbreak{}accuracy\_\allowbreak{}of\_\allowbreak{}error\_\allowbreak{}budgets} gives the formal counterpart of this step. Named proof applications include \texttt{first\allowbreak{}Order\_\allowbreak{}accuracy\_\allowbreak{}budget}, \texttt{Canonical\allowbreak{}Coupling.\allowbreak{}klmc\_\allowbreak{}wasserstein\_\allowbreak{}bound}; the adjacent proof shows their hypotheses and order of use.
\end{leanbox}


Earlier results used here: \cite[Section~3.17.3]{LangevinLean2026ULASC}; \cite[Section~3.17.4]{LangevinLean2026ULASC}; \cite[Section~3.26.3]{LangevinLean2026ULASC}.

\section{\texorpdfstring{Choose the paper step and round the count}{Choose the paper step and round the count}}


For positive dimension choose \(h=\min\{m/(4\gamma M),\,0.94\varepsilon m/(M\sqrt{2p})\}\). Both terms are positive. The first enforces stability and the second enforces the discretization-error bound by multiplication through positive denominators. Take the corresponding nonnegative step length, use the rounded logarithmic count, and apply the accuracy theorem. At \(\gamma=\sqrt{m+M}\), this is precisely the manuscript’s minimum-step prescription. The explicit reciprocal dependence displays the \(\sqrt p/\varepsilon\) iteration scaling up to the logarithm.



The chosen step is the minimum of a stability limit and an accuracy limit. Taking that minimum guarantees both requirements simultaneously. Its reciprocal makes explicit the tradeoff between small requested error, dimension-dependent stationary velocity fluctuations, and the curvature-dependent contraction rate.

\begin{leanbox}{Lean: \texttt{klmcAccuracyStep}, \texttt{klmcAccuracyStep\_\allowbreak{}spec}, \texttt{klmc\_\allowbreak{}accuracy\_\allowbreak{}explicit}}
\leansource{\texttt{klmcAccuracyStep} \quad KineticLangevin/\allowbreak{}Accuracy.lean, lines 41–42}
\begin{leancode}
\leanline{0}{def\ klmcAccuracyStep\ (m\ M\ γ\ p\ ε\ :\ ℝ)\ :\ ℝ\ :=}
\leanline{2}{min\ (m/(4*γ*M))\ ((47/50:ℝ)*ε*m/(M*Real.sqrt\ (2*p)))}
\end{leancode}
\leansource{\texttt{klmcAccuracyStep\_\allowbreak{}spec} \quad KineticLangevin/\allowbreak{}Accuracy.lean, lines 45–60}
\begin{leancode}
\leanline{0}{theorem\ klmcAccuracyStep\_spec\ (m\ M\ γ\ p\ ε\ :\ ℝ)\ (hm\ :\ 0<m)\ (hM\ :\ 0<M)}
\leanline{4}{(hγ\ :\ 0<γ)\ (hp\ :\ 0<p)\ (hε\ :\ 0<ε)\ :}
\leanline{4}{0<klmcAccuracyStep\ m\ M\ γ\ p\ ε\ ∧}
\leanline{4}{4*γ*M*klmcAccuracyStep\ m\ M\ γ\ p\ ε≤m\ ∧}
\leanline{4}{M*klmcAccuracyStep\ m\ M\ γ\ p\ ε*Real.sqrt\ (2*p)/m≤(47/50:ℝ)*ε}
\end{leancode}
\leansource{\texttt{klmc\_\allowbreak{}accuracy\_\allowbreak{}explicit} \quad KineticLangevin/\allowbreak{}Accuracy.lean, lines 63–77}
\begin{leancode}
\leanline{0}{theorem\ klmc\_accuracy\_explicit\ \{d\ :\ ℕ\}\ (hd\ :\ 0<d)\ \{μ\ :\ Measure\ (CanonicalCoupling.H\ d)\}}
\leanline{4}{[IsProbabilityMeasure\ μ]\ (f\ :\ CanonicalCoupling.H\ d\ →\ ℝ)\ (m\ :\ ℝ)\ (M\ γ\ :\ ℝ≥0)}
\leanline{4}{(hf\ :\ Differentiable\ ℝ\ f)\ (hsc\ :\ StrongConvexOn\ univ\ m\ f)\ (hLip\ :\ LipschitzWith\ M\ (gradient\ f))}
\leanline{4}{(hmM\ :\ m≤(M:ℝ))\ (hm\ :\ 0<m)\ (hγ\ :\ 0<(γ:ℝ))\ (hfriction\ :\ m+(M:ℝ)≤(γ:ℝ)\textasciicircum{}2)}
\leanline{4}{(hfinite\ :\ ULA.wasserstein2\ μ\ (ULA.gibbsMeasure\ f)≠⊤)\ (ε\ :\ ℝ)\ (hε\ :\ 0<ε)\ :}
\leanline{4}{let\ hs\ :=\ klmcAccuracyStep\_spec\ m\ M\ γ\ d\ ε\ hm\ (hm.trans\_le\ hmM)\ hγ\ (by\ exact\_mod\_cast\ hd)\ hε}
\leanline{4}{let\ h\ :\ ℝ≥0\ :=\ ⟨klmcAccuracyStep\ m\ M\ γ\ d\ ε,hs.1.le⟩}
\leanline{4}{let\ n\ :=\ firstOrderIterationCount\ (3*m*(h:ℝ)/(4*(γ:ℝ)))\ (ULA.wasserstein2Real\ μ\ (ULA.gibbsMeasure\ f))\ ε}
\leanline{4}{ULA.wasserstein2Real\ (CanonicalCoupling.klmcLaw\ f\ γ\ h\ μ\ n)\ (ULA.gibbsMeasure\ f)≤ε}
\end{leancode}
\leanlabel{How the Lean proof follows the argument}
\texttt{klmc\allowbreak{}Accuracy\allowbreak{}Step} names the mathematical object defined above. Its defining expression appears in the adjacent statement; subsequent identities may unfold this definition directly.




\texttt{klmc\allowbreak{}Accuracy\allowbreak{}Step\_\allowbreak{}spec} gives the formal counterpart of this step. The main intermediate facts are \texttt{hs} in that order. The final step is the algebraic normalization or ordered arithmetic displayed in the proof.




\texttt{klmc\_\allowbreak{}accuracy\_\allowbreak{}explicit} gives the formal counterpart of this step. The main intermediate facts are \texttt{hs}, \texttt{h} in that order. Named proof applications include \texttt{klmc\_\allowbreak{}accuracy\_\allowbreak{}of\_\allowbreak{}error\_\allowbreak{}budgets}; the adjacent proof shows their hypotheses and order of use.
\end{leanbox}


Earlier results used here: \cite[Section~3.17.3]{LangevinLean2026ULASC}; \cite[Section~3.17.4]{LangevinLean2026ULASC}; \cite[Section~3.26.3]{LangevinLean2026ULASC}.

\begin{thebibliography}{ULA-LSI}
\bibitem[ULA-SC]{LangevinLean2026ULASC}
LangevinLean Project. \emph{Unadjusted Langevin under Strong Convexity}. LL-01, 2026. Draft report, content revision d50c0f64815b. \url{https://odakyildiz.github.io/LangevinLean/docs/pdf/UnadjustedLangevin/FullProof.pdf}.
\end{thebibliography}

\end{document}
